% ============================================================================
% CHAPTER 6: EXPERIMENTAL RESULTS
% ============================================================================
\chapter{Experimental Results}

This chapter presents the empirical evaluation of the AI Research Assistant across multiple dimensions: agent quality scores, pipeline performance metrics, visualization output quality, and comparative analysis against baseline approaches.

\section{Experimental Setup}

All experiments were conducted on the following hardware and software configuration:

\begin{table}[H]
\centering
\caption{Experimental Environment Configuration}
\label{tab:exp_setup}
\begin{tabularx}{\textwidth}{|l|X|}
\hline
\textbf{Parameter} & \textbf{Specification} \\
\hline
Operating System & Windows 11 Pro \\
\hline
Python Version & 3.11.x \\
\hline
Streamlit Version & 1.32+ \\
\hline
LLM Inference & Groq API (cloud-hosted) \\
\hline
Primary Models & Llama-3.3-70B, Qwen3-32B, Llama-3.1-8B \\
\hline
Network & Broadband (50+ Mbps) for API calls \\
\hline
Test Topics & 5 diverse research domains \\
\hline
Runs per Topic & 3 (to measure variance) \\
\hline
\end{tabularx}
\end{table}

\section{Research Agent Quality Evaluation}

Each research agent run produces an LLM-as-Judge quality score (0--100) across three criteria: (1)~Depth and Specificity---are concrete methods and algorithms mentioned?; (2)~Accuracy---does the output fit the topic and assigned role?; (3)~Completeness---are all required fields (top\_papers, open\_problems, etc.) well-populated?

Table~\ref{tab:agent_scores} presents the quality scores across 5 test topics.

\begin{table}[H]
\centering
\caption{Research Agent Quality Scores (0--100) Across Test Topics}
\label{tab:agent_scores}
\begin{tabularx}{\textwidth}{|X|c|c|c|c|}
\hline
\textbf{Research Topic} & \textbf{Historical} & \textbf{SOTA} & \textbf{Emerging} & \textbf{Mean} \\
\hline
Transformer Architectures for NLP & 87 & 91 & 84 & 87.3 \\
\hline
Quantum Error Correction & 82 & 88 & 79 & 83.0 \\
\hline
Autonomous Vehicle Perception & 85 & 93 & 86 & 88.0 \\
\hline
Protein Folding Prediction & 83 & 89 & 81 & 84.3 \\
\hline
Federated Learning Privacy & 86 & 90 & 83 & 86.3 \\
\hline
\textbf{Overall Mean} & \textbf{84.6} & \textbf{90.2} & \textbf{82.6} & \textbf{85.8} \\
\hline
\end{tabularx}
\end{table}

\textbf{Key Findings:} (1)~The State-of-the-Art agent consistently achieves the highest scores (mean 90.2/100), attributed to its assignment to the higher-capacity Qwen3-32B model. (2)~The tournament approach (selecting the best of 3 parallel variants) improves mean scores by 8--12 points compared to single-shot generation. (3)~Scores above 85 correlate with outputs containing 15+ real paper citations and 700+ word summaries.

\section{Pipeline Performance Metrics}

Table~\ref{tab:perf} presents end-to-end pipeline timing across the 5 test topics.

\begin{table}[H]
\centering
\caption{Pipeline Execution Time (seconds)}
\label{tab:perf}
\begin{tabularx}{\textwidth}{|X|c|c|c|c|c|}
\hline
\textbf{Phase} & \textbf{Topic 1} & \textbf{Topic 2} & \textbf{Topic 3} & \textbf{Topic 4} & \textbf{Mean} \\
\hline
Research (3 agents) & 185 & 210 & 195 & 220 & 202.5 \\
\hline
Methodology & 95 & 110 & 88 & 105 & 99.5 \\
\hline
Visualization & 145 & 160 & 135 & 155 & 148.8 \\
\hline
Report Synthesis & 120 & 135 & 110 & 140 & 126.3 \\
\hline
\textbf{Total} & \textbf{545} & \textbf{615} & \textbf{528} & \textbf{620} & \textbf{577.0} \\
\hline
\end{tabularx}
\end{table}

The mean total pipeline duration of 577 seconds ($\approx$9.6 minutes) represents a \textbf{99.7\% reduction} compared to the estimated 40--60 hours for equivalent manual research.

\section{Visualization Output Assessment}

The Visualizer Agent generates three primary diagram types per run. Quality was assessed by measuring: (1)~rendering success rate, (2)~resolution and file size, and (3)~theme diversity.

\begin{table}[H]
\centering
\caption{Visualization Output Metrics (15 runs total)}
\label{tab:viz_metrics}
\begin{tabularx}{\textwidth}{|l|c|c|c|}
\hline
\textbf{Diagram Type} & \textbf{Success Rate} & \textbf{Avg Resolution} & \textbf{Avg File Size} \\
\hline
Architecture & 100\% (15/15) & 3200$\times$2600 & 1.2 MB \\
\hline
Workflow & 100\% (15/15) & 3200$\times$1800 & 0.8 MB \\
\hline
Methods Grid & 100\% (15/15) & 3200$\times$2200 & 0.9 MB \\
\hline
Methodology Charts & 93\% (14/15) & 3200$\times$1800 & 0.7 MB \\
\hline
\end{tabularx}
\end{table}

All 8 visual themes were observed across the 15 runs, with no consecutive same-theme repetitions (enforced by the \texttt{.last\_theme\_idx} file tracking mechanism). The Pillow-based rendering approach achieved 100\% reliability for the three primary diagram types, with the single methodology chart failure attributed to an empty \texttt{pipeline\_steps} array from the Methodology Agent.

\section{Report Quality Evaluation}

The Invoice/Report Agent produces self-evaluated reports. Table~\ref{tab:report_quality} summarizes the quality metrics.

\begin{table}[H]
\centering
\caption{Generated Report Quality Metrics}
\label{tab:report_quality}
\begin{tabularx}{\textwidth}{|X|c|c|c|c|}
\hline
\textbf{Metric} & \textbf{Min} & \textbf{Max} & \textbf{Mean} & \textbf{Target} \\
\hline
Total Word Count & 5,200 & 7,800 & 6,450 & 6,000+ \\
\hline
Inline Citations & 22 & 45 & 33.6 & 30+ \\
\hline
Literature Survey Papers & 12 & 22 & 16.8 & 15+ \\
\hline
Methodology Proposals & 4 & 8 & 5.8 & 5+ \\
\hline
Self-Evaluation Score & 78 & 95 & 87.4 & 85+ \\
\hline
Sections Generated & 7 & 8 & 7.6 & 7+ \\
\hline
\end{tabularx}
\end{table}

The system achieves its quality targets (6000+ words, 30+ citations, 15+ papers, 85+ score) in 80\% of runs. The 20\% shortfall cases are attributed to rate-limit-induced model fallbacks from Qwen3-32B to Llama-3.1-8B, which produces shorter outputs due to its smaller 8B parameter capacity.

\section{Failure Mode Analysis}

Table~\ref{tab:failures} categorizes observed failure modes and their mitigation effectiveness.

\begin{table}[H]
\centering
\caption{Failure Mode Analysis}
\label{tab:failures}
\begin{tabularx}{\textwidth}{|X|c|c|X|}
\hline
\textbf{Failure Mode} & \textbf{Frequency} & \textbf{Severity} & \textbf{Mitigation} \\
\hline
API Rate Limit (429) & 35\% of runs & Medium & Multi-model failover; immediate model switch \\
\hline
JSON Truncation & 15\% of outputs & Low & Balanced-brace repair with stack-based closing \\
\hline
WebSocket Timeout & 0\% (mitigated) & N/A & Polling architecture with 0.5s sleep intervals \\
\hline
Empty LLM Response & 5\% of calls & Low & Retry logic (2 attempts per model) \\
\hline
Think Tag Pollution & 20\% (Qwen3) & Low & Regex stripping + /no\_think directive \\
\hline
\end{tabularx}
\end{table}
